% !TEX program = xelatex
% 上海电机学院 毕业设计（论文）答辩示例。编译：latexmk slides.tex
\documentclass{sdjuslides}
\usetikzlibrary{positioning}
% 合成回放数据。每轮向两种规则输入相同的 100 个正常窗口。
% 每个命令的值为该轮触发告警的正常窗口数。
\newcommand\windowsPerTrial{100}
\newcommand\singleTrialA{12}
\newcommand\singleTrialB{10}
\newcommand\singleTrialC{11}
\newcommand\doubleTrialA{4}
\newcommand\doubleTrialB{5}
\newcommand\doubleTrialC{3}


% 封面信息、术语、单位、数据和结论应与论文定稿一致。
% 示例的身份和数据均为占位或合成内容。答辩时替换为本人信息和可核对的成果。
\title[电机振动异常监测]{基于窗口判定的电机振动异常监测}
\titleen{Window-Based Anomaly Monitoring of Motor Vibration}
\author{测试学生甲}
\studentid{XXXXXXXXXX}
\major{电气工程及其自动化}
\school{电气学院}
\supervisor{XXX}
\companysupervisor{XXX} % 没有企业指导教师时可省略，封面该栏留空
\defensedate{20XX 年 XX 月 XX 日}

\begin{document}

\makecover
\sdjutableofcontents

\section{研究背景与问题}
\begin{frame}{短时振动尖峰会触发误报}
  电机在正常运行时也可能出现瞬时振动尖峰，单个窗口越限就告警会打断巡检。
  \medskip
  \begin{itemize}
    \item 研究问题：要求连续越限，能否减少正常窗口的误报？
    \item 完成工作：实现两种告警规则，并用同一批窗口回放对比。
  \end{itemize}
\end{frame}

\section{方案与设计}
\begin{frame}{同一评分，两种告警规则}
  \begin{figure}
    \centering
    \begin{tikzpicture}[
    workflow step/.style={draw=sdjuaccent,fill=sdjusurface,rounded corners=2pt,
      minimum height=3em,inner xsep=.8em},
    flow/.style={->,semithick,color=sdjuaccent},
  ]
  \node[workflow step] (replay) {窗口回放};
  \node[workflow step,right=of replay] (score) {振动评分};
  \node[workflow step,right=of score] (judge) {连续判定};
  \node[workflow step,right=of judge] (log) {记录告警};
  \draw[flow] (replay) -- (score);
  \draw[flow] (score) -- (judge);
  \draw[flow] (judge) -- (log);
\end{tikzpicture}

    \caption{窗口回放与告警流程}
    \sdjusource{本示例绘制，\path{slides/figures/workflow.tex}}
  \end{figure}
  \begin{columns}[T,onlytextwidth]
    \begin{column}{.47\textwidth}
      \textbf{单窗规则}\par
      当前窗口评分越限即告警。
    \end{column}
    \begin{column}{.47\textwidth}
      \textbf{双窗规则}\par
      当前窗口和上一窗口均越限才告警。
    \end{column}
  \end{columns}
\end{frame}

\section{实现与实验}
\begin{frame}[fragile]{逐窗判定是否告警}
  两种规则共用窗口评分和阈值。每轮开头将上一窗口的越限状态设为假。
  \medskip
  \begin{lstlisting}[language=Python]
prev_above = False
for score in scores:
    above = score > threshold
    single_alarm = above
    double_alarm = prev_above and above
    prev_above = above
  \end{lstlisting}
\end{frame}

\begin{frame}{用正常窗口回放评估误报}
  回放三轮，每轮 \windowsPerTrial{} 个正常窗口。两种规则使用完全相同的输入。
  \[
    \mathrm{FPR} = \frac{N_{\mathrm{FP}}}{N_{\mathrm{normal}}} \times 100\%
  \]
  分子是触发告警的正常窗口数。分母是参与回放的正常窗口总数。
\end{frame}

\section{结果与数据}
\begin{frame}{三轮回放均观察到误报减少}
  \begin{table}
    \centering
    \caption{合成回放的误报窗口数}
    \begin{tabular}{lrrr}
      \toprule
      轮次 & 正常窗口 & 单窗误报 & 双窗误报 \\
      \midrule
      第一轮 & \windowsPerTrial & \singleTrialA & \doubleTrialA \\
      第二轮 & \windowsPerTrial & \singleTrialB & \doubleTrialB \\
      第三轮 & \windowsPerTrial & \singleTrialC & \doubleTrialC \\
      \bottomrule
    \end{tabular}
    \sdjusource{合成数据，\path{slides/data/replay-counts.tex}}
  \end{table}
\end{frame}

\begin{frame}{合成回放的误报率由 11\% 降至 4\%}
  \pgfmathsetmacro{\singleRate}{100*(\singleTrialA+\singleTrialB+\singleTrialC)/(3*\windowsPerTrial)}
  \pgfmathsetmacro{\doubleRate}{100*(\doubleTrialA+\doubleTrialB+\doubleTrialC)/(3*\windowsPerTrial)}
  \begin{figure}
    \centering
    \begin{tikzpicture}[x=.05\textwidth,y=1.45em]
      \foreach \tick in {0,4,8,12} {
        \draw[sdjutext!18] (\tick,.58) -- (\tick,2.45);
        \node[below,text=sdjumuted] at (\tick,.58) {\tick\%};
      }
      \fill[sdjumuted] (0,1.73) rectangle (\singleRate,2.27);
      \fill[sdjuaccent] (0,.73) rectangle (\doubleRate,1.27);
      \node[left] at (0,2) {单窗};
      \node[left] at (0,1) {双窗};
      \node[right] at (\singleRate,2) {\pgfmathprintnumber{\singleRate}\%};
      \node[right] at (\doubleRate,1) {\pgfmathprintnumber{\doubleRate}\%};
    \end{tikzpicture}
    \caption{两种规则的合计误报率}
    \sdjusource{由 \path{slides/data/replay-counts.tex} 的三轮计数计算}
  \end{figure}
\end{frame}

\section{结论与不足}
\begin{frame}{双窗规则降低误报，检出能力待验证}
  \begin{itemize}
    \item 三轮合成回放中，双窗规则的误报率低了 7 个百分点。
    \item 回放仅含正常窗口，无法估计漏报率。持续越限时，双窗规则晚一个窗口告警。
    \item 后续用带故障标签的实测数据验证检出率和告警延迟。
  \end{itemize}
\end{frame}

\sdjuthanks{感谢指导教师和实验室同学的帮助，敬请各位老师批评指正。}

\end{document}
